\section{Déploiement, sécurité et exploitation}

\subsection{Déploiement AWS}
\diagramme{09_deploiement_aws}{N\oe{}uds, services et artefacts déployés}{fig:9}

\subsubsection{Principes}

\begin{itemize}
\item Une région unique est paramétrée dans l'infrastructure comme code.
\item Les ressources portent des noms et étiquettes permettant de les retrouver et de les supprimer.
\item Le frontend ne contient aucun secret AWS.
\item L'instance EC2 utilise un rôle IAM limité à SQS, S3, SES et CloudWatch.
\item Le bucket de PDF refuse l'accès public direct. L'accès est accordé de manière temporaire après réservation.
\item Les files possèdent une politique de reprise et une file d'échecs contrôlée.
\end{itemize}

\subsubsection{Infrastructure comme code}

La description de l'infrastructure doit permettre la création et la suppression reproductibles de l'instance EC2, de la base RDS PostgreSQL, des files, des buckets, des rôles, des journaux et des alarmes. Les paramètres dépendant de l'environnement sont séparés du code. La commande exacte sera documentée dans le README.


\subsection{Sécurité}
\subsubsection{Moindre privilège}

\begingroup\small\setlength{\tabcolsep}{3pt}\renewcommand{\arraystretch}{1.25}
\begin{longtable}{|P{0.22\textwidth}|P{0.36\textwidth}|P{0.35\textwidth}|}
\caption{Permissions des composants}\\
\hline
Composant & Autorisations nécessaires & Autorisations exclues \\ \hline
\endfirsthead
\hline
Composant & Autorisations nécessaires & Autorisations exclues \\ \hline
\endhead
\endfoot

Backend Node.js & Accéder à PostgreSQL, publier les travaux dans SQS et écrire les journaux CloudWatch. & Administrer le compte AWS ou rendre le bucket public. \\ \hline

Service de téléchargement & Lire et mettre à jour un lien dans PostgreSQL, créer une réservation et placer un travail dans SQS. & Modifier les opérations d'édition ou les crédits. \\ \hline

Worker PDF & Lire une version, écrire le PDF correspondant dans S3 et mettre à jour le statut. & Décrémenter le quota. \\ \hline

Worker de notification & Lire un événement de dernier exemplaire, appeler SES et marquer la notification. & Lire le contenu complet des documents. \\ \hline

Administration & Lire les vues opérationnelles. & Modifier silencieusement les quotas ou l'historique. \\ \hline

\end{longtable}\endgroup

\subsubsection{Mesures applicatives}

\begin{itemize}
\item validation de la syntaxe et de la taille des entrées.
\item jetons de partage imprévisibles.
\item limitation des requêtes par identité déclarée, session et signal réseau pertinent.
\item échappement du contenu lors de l'affichage et de la génération PDF.
\item absence de données sensibles dans les journaux.
\item expiration rapide des URL S3 présignées \cite{awspresigned}.
\item chiffrement des communications et chiffrement au repos fourni par les services utilisés.
\end{itemize}

\subsubsection{Limite de l'identification}

L'identification par nom et e-mail est déclarative. Un utilisateur peut donc déclarer l'adresse d'un tiers. Les autorisations et les blocages portent sur l'identité applicative présentée. Ils ne garantissent pas l'identité de la personne physique. Cette limite détermine la portée réelle du contrôle d'accès et du dispositif anti-abus.


\subsection{Performance, robustesse et observabilité}
\subsubsection{Traitement du pic}

\diagramme{10_activite_pic}{Réponse rapide et arrêt des travaux inutiles}{fig:10}

Le chemin synchrone Express décide rapidement si un exemplaire peut être réservé. La génération, plus coûteuse, est découplée par SQS. Une demande refusée n'entre jamais dans la file. Le processus WebSocket Yjs est séparé du worker PDF afin que le trafic public ne consomme pas toute la capacité collaborative. À l'épuisement du lien, les réservations déjà acceptées poursuivent leur traitement. Les nouvelles demandes reçoivent un état épuisé. L'arrêt des travaux inutiles ne doit pas interrompre une livraison autorisée.


\subsubsection{Modes de défaillance}

\begingroup\small\setlength{\tabcolsep}{3pt}\renewcommand{\arraystretch}{1.25}
\begin{longtable}{|P{0.22\textwidth}|P{0.36\textwidth}|P{0.35\textwidth}|}
\caption{Gestion des défaillances}\\
\hline
Défaillance & Réaction prévue & Garantie conservée \\ \hline
\endfirsthead
\hline
Défaillance & Réaction prévue & Garantie conservée \\ \hline
\endhead
\endfoot

Message PDF répété & Réutiliser la génération identifiée par réservation/version. & Pas de double consommation. \\ \hline

Client collaboratif déconnecté & Recharger l'instantané puis les opérations manquantes. & Convergence finale. \\ \hline

E-mail temporairement indisponible & Nouvelle tentative idempotente. & Événement unique et livraison externe suivie, avec un risque résiduel de doublon. \\ \hline

PDF échoué & État ÉCHEC visible et reprise bornée. & Pas de fichier incohérent. \\ \hline

Lien déjà épuisé & Réponse immédiate sans travail asynchrone. & Aucun service inexistant promis. \\ \hline

\end{longtable}\endgroup

\subsubsection{Observabilité}

\begingroup\small\setlength{\tabcolsep}{3pt}\renewcommand{\arraystretch}{1.25}
\begin{longtable}{|P{0.26\textwidth}|P{0.69\textwidth}|}
\caption{Indicateurs d’exploitation}\\
\hline
Indicateur & Utilité \\ \hline
\endfirsthead
\hline
Indicateur & Utilité \\ \hline
\endhead
\endfoot

Réservations acceptées/refusées par lien & Prouver l'exactitude et diagnostiquer les pics. \\ \hline

Latence p50/p95/p99 de réservation & Mesurer la qualité de réponse publique. \\ \hline

Profondeur et âge maximal de la file PDF & Détecter une saturation. \\ \hline

Connexions et opérations WebSocket & Observer les documents actifs. \\ \hline

Retards de propagation & Vérifier « presque aussitôt ». \\ \hline

Refus de crédit et blocages & Suivre les limites et les abus. \\ \hline

Erreurs et reprises de génération & Contrôler la robustesse. \\ \hline

Notification du dernier exemplaire & Vérifier l'unicité de l'événement. \\ \hline

\end{longtable}\endgroup
